\section{Metodologia}
\label{sec:metodologia}

O estudo foi organizado em seis etapas: construção da base agrícola, consolidação
da base meteorológica, preparação das séries diárias para o DSSAT, calibração da
produtividade simulada, formulação do ambiente de aprendizado por reforço e
treinamento do agente PPO. O fluxo foi implementado em Python e a simulação das
safras foi realizada com o DSSAT-CSM por meio do pacote
\texttt{Py\_DSSATTools}. A calibração e a avaliação foram feitas por divisão
temporal, de modo que os anos usados para teste não participassem do ajuste do
modelo.

\subsection{Base de produtividade agrícola}
\label{subsec:base_produtividade}

Os dados agrícolas foram obtidos na Tabela 1612 do Sistema IBGE de Recuperação
Automática (SIDRA), disponível em
\url{https://sidra.ibge.gov.br/tabela/1612}. A consulta foi feita para lavouras
temporárias em escala municipal, mantendo os municípios disponíveis na base
oficial. O arquivo exportado continha oito tabelas, cada uma com 5.563 linhas e
30 colunas, totalizando 1.335.120 valores antes da etapa de recorte e
normalização.

Foram mantidas as oito variáveis disponíveis na consulta:
\begin{enumerate}
    \item área plantada, em hectares;
    \item área plantada como percentual do total geral;
    \item área colhida, em hectares;
    \item área colhida como percentual do total geral;
    \item quantidade produzida, em toneladas;
    \item rendimento médio da produção, em kg ha$^{-1}$;
    \item valor da produção, na unidade monetária indicada pelo SIDRA em cada ano;
    \item valor da produção como percentual do total geral.
\end{enumerate}

A variável de valor da produção foi preservada com a unidade original informada
pelo SIDRA, sem conversão entre moedas. Essa variável não foi usada como alvo do
modelo, mas foi mantida na base final para rastreabilidade. A cultura de
interesse para a simulação foi soja em grão. A série agrícola da consulta
abrangeu o período de 1974 a 2024. Para este experimento, foram extraídos os
registros dos municípios de Castro (PR) e Ponta Grossa (PR), preservando ano,
município, produto, variável, valor original e valor numérico quando disponível.
Valores ausentes ou marcados na fonte como sem informação foram mantidos como
ausentes, sem imputação.

A base agrícola consolidada resultou em 8.640 registros em formato longo. Para a
calibração e avaliação da produtividade simulada, foi usado o rendimento médio
da produção de soja em grão em Castro, expresso em kg ha$^{-1}$.

\subsection{Base meteorológica}
\label{subsec:base_meteorologica}

Os dados meteorológicos foram obtidos no portal de dados históricos do Instituto
Nacional de Meteorologia (INMET), disponível em
\url{https://portal.inmet.gov.br/dadoshistoricos}. Foi usada a estação
automática A819, localizada em Castro, Paraná. As séries horárias incluíram
precipitação, temperatura do ar, temperatura máxima e mínima na hora anterior,
radiação global, umidade relativa, pressão atmosférica, direção e velocidade do
vento.

Foram lidos os arquivos anuais disponíveis entre 2006 e 2026. O arquivo
\texttt{dados\_A819\_H\_2006-07-08\_2026-01-01.csv} foi usado como fonte
complementar quando havia lacunas no histórico anual. A consolidação manteve os
registros horários dos arquivos anuais e adicionou apenas os horários ausentes.
Esse procedimento incluiu 8.784 linhas da fonte complementar: 24 horas de 8 de
julho de 2006 e 8.760 horas referentes ao ano de 2022. A base meteorológica
final ficou com 174.432 registros horários, sem duplicidade em
\texttt{data\_hora}, no intervalo de 8 de julho de 2006 a 31 de maio de 2026.

Antes da agregação diária, valores sentinela usados nos arquivos meteorológicos
para indicar ausência de observação foram convertidos para ausentes. Valores
menores ou iguais a $-90$ foram tratados como ausentes nas variáveis numéricas.
Também foram removidos valores fisicamente inconsistentes, como precipitação
negativa, temperaturas fora do intervalo de $-20$ a $50\,^\circ$C e radiação
global horária negativa ou superior a 6000 kJ m$^{-2}$.

A agregação diária foi feita a partir das observações horárias válidas. A
precipitação diária foi calculada pela soma horária; a temperatura máxima e
mínima diárias foram obtidas pelos extremos horários; a temperatura média foi
calculada pela média das observações horárias; a radiação global diária foi
calculada pela soma horária e convertida de kJ m$^{-2}$ para MJ m$^{-2}$.
Para reduzir o efeito de dias incompletos, dias com menos de 18 observações de
precipitação tiveram a chuva diária tratada como ausente; dias com menos de 18
observações de temperatura tiveram as temperaturas tratadas como ausentes; dias
com menos de 6 observações de radiação tiveram a radiação tratada como ausente.

As falhas de precipitação foram preenchidas pela média climatológica do dia do
ano, com média mensal como alternativa quando necessário. A radiação ausente ou
inferior a 0,5 MJ m$^{-2}$ dia$^{-1}$ foi estimada pelo método de Hargreaves, a
partir da amplitude térmica diária e da radiação extraterrestre calculada para a
latitude da estação. Essa estimativa foi usada apenas quando a observação do
INMET não estava disponível em condição adequada para alimentar o DSSAT.

\subsection{Divisão temporal}
\label{subsec:divisao_temporal}

Cada safra foi definida pelo ano de plantio e pela janela meteorológica de
setembro do ano corrente a abril do ano seguinte. A divisão temporal usada no
experimento foi:
\begin{itemize}
    \item treinamento: 2006--2017;
    \item validação: 2018--2021;
    \item teste: 2022--2024 para comparação com produtividade observada;
    \item simulação climática adicional: 2025, sem produtividade SIDRA observada
    na base disponível no momento da análise.
\end{itemize}

Essa separação foi adotada para evitar que uma mesma safra participasse do
ajuste e da avaliação final. Os anos de validação foram usados para seleção da
correção estatística e para acompanhamento do treinamento PPO. Os anos de teste
foram usados apenas na avaliação final da produtividade.

\subsection{Simulação da cultura no DSSAT}
\label{subsec:dssat}

O DSSAT (\textit{Decision Support System for Agrotechnology Transfer}) foi usado
como ambiente mecanístico de simulação da cultura. Diferentemente de uma
regressão ajustada diretamente sobre os dados históricos, o DSSAT simula o ciclo
da soja a partir de processos diários de desenvolvimento fenológico, crescimento,
balanço hídrico do solo, acúmulo de biomassa e formação de grãos. Assim, cada
combinação de clima, solo, cultivar, data de plantio e irrigação gera uma safra
simulada.

Foi usado o módulo de soja do DSSAT, \texttt{CRGRO-Soybean}. A integração com
Python foi feita com \texttt{Py\_DSSATTools}, que cria os arquivos de entrada,
executa o DSSAT-CSM e lê os resultados de saída. Para cada simulação, o código
gera um tratamento contendo:
\begin{enumerate}
    \item série meteorológica diária da estação A819;
    \item perfil de solo no formato DSSAT;
    \item cultivar de soja;
    \item data de plantio;
    \item população, espaçamento e profundidade de semeadura;
    \item eventos de irrigação;
    \item controles de simulação e arquivos de saída.
\end{enumerate}

A estação foi parametrizada com latitude $-24{,}77999999$, longitude
$-50{,}04638888$ e altitude de 1.016 m. O arquivo meteorológico diário fornecido
ao DSSAT continha radiação solar, temperatura máxima, temperatura mínima e
precipitação. A cultivar configurada foi \texttt{990005}, disponível no arquivo
\texttt{SBGRO048.CUL}. A população foi fixada em 30 plantas m$^{-2}$, com
espaçamento entre linhas de 45 cm e profundidade de semeadura de 5 cm.

Na versão computacional usada neste estudo, o solo foi representado por um perfil
argiloso padrão para Castro, com cinco camadas até 150 cm. O perfil contém limite
inferior de água, capacidade de campo, saturação, densidade do solo, carbono
orgânico, textura e fator de crescimento radicular. Esse perfil permite executar
o fluxo completo com DSSAT real. Para aplicação local em campo, recomenda-se
substituir esse perfil por um \texttt{SOIL.SOL} medido ou calibrado para a área
experimental.

\subsection{Calibração da produtividade simulada}
\label{subsec:calibracao_dssat}

Como o DSSAT depende de solo, cultivar e parâmetros de manejo que nem sempre
estão disponíveis em escala municipal, foi adicionada uma etapa de correção
estatística da produtividade simulada. O objetivo dessa etapa foi alinhar a
produtividade simulada pelo DSSAT ao rendimento médio municipal observado no
SIDRA, mantendo o DSSAT como gerador da resposta agronômica de cada manejo.

Para cada ano dos conjuntos de treinamento, validação e teste, o DSSAT foi
executado em cinco datas de plantio: 25 de setembro, 5 de outubro, 15 de
outubro, 25 de outubro e 5 de novembro. Nessas simulações de calibração, foi
usada uma política fixa de irrigação com gatilho de déficit hídrico igual a
45 mm, lâmina de 20 mm por evento e limite sazonal de 120 mm. Para cada
simulação foram registrados produtividade bruta do DSSAT, precipitação,
irrigação, dia juliano de plantio, temperatura média e radiação média da safra.

A correção foi ajustada por regressão ridge. O vetor de entrada incluiu ano,
produtividade bruta do DSSAT, precipitação, irrigação, dia juliano de plantio,
temperatura média, radiação média e termos quadráticos e de interação derivados
dessas variáveis. O modelo geral pode ser escrito como:
\begin{equation}
    \hat{Y}_{\mathrm{corr}} =
    \beta_0 + \sum_{j=1}^{p} \beta_j x_j,
\end{equation}
em que $\hat{Y}_{\mathrm{corr}}$ é a produtividade corrigida, $x_j$ são as
variáveis derivadas da simulação DSSAT e $\beta_j$ são os coeficientes estimados.
Os coeficientes foram obtidos pela minimização:
\begin{equation}
    \min_{\boldsymbol{\beta}}
    \sum_{i=1}^{n}
    \left(Y_i - \hat{Y}_{i,\mathrm{corr}}\right)^2
    + \alpha \sum_{j=1}^{p}\beta_j^2,
\end{equation}
em que $Y_i$ é a produtividade observada no SIDRA e $\alpha$ é o parâmetro de
regularização.

A busca de $\alpha$ foi feita no conjunto
\[
\{0{,}05,\ 0{,}10,\ 0{,}50,\ 1,\ 2{,}5,\ 5,\ 10,\ 25,\ 50,\ 100,\ 200,\ 300,\ 500\}.
\]
O ajuste dos coeficientes foi feito com os anos de treinamento. Para cada valor
de $\alpha$, o modelo foi aplicado às simulações da validação, as previsões foram
agregadas por ano, e o erro médio de validação foi calculado. Em seguida, foi
aplicado um ajuste de intercepto igual ao viés médio observado na validação. O
parâmetro final selecionado foi $\alpha = 200$, por apresentar o menor erro
absoluto médio após o ajuste de intercepto na validação. O modelo final de
correção foi salvo em \texttt{outputs/calibration/yield\_correction.json} e
passou a ser aplicado automaticamente após cada simulação DSSAT.

\subsection{Ambiente de aprendizado por reforço}
\label{subsec:ambiente_rl}

O ambiente foi implementado com a interface \texttt{Gymnasium}. Cada episódio
representa uma safra completa. No início do episódio, um ano meteorológico é
selecionado dentro do conjunto em uso. O agente observa estatísticas climáticas
da safra e escolhe, em uma única ação, a data de plantio e os parâmetros de
irrigação. Em seguida, o ambiente constrói o tratamento DSSAT, executa a
simulação e retorna produtividade, água aplicada e recompensa.

O estado observado pelo agente foi:
\begin{equation}
    s =
    \left[
    \frac{P}{1800},
    \frac{\bar{T}}{35},
    \frac{T_{\max}}{45},
    \frac{T_{\min}}{20},
    \frac{\bar{R}_{s}}{30},
    \frac{n_{\mathrm{dias}}}{366},
    0,
    0,
    1
    \right],
\end{equation}
em que $P$ é a precipitação acumulada da safra, $\bar{T}$ é a temperatura média,
$T_{\max}$ e $T_{\min}$ são os extremos térmicos, $\bar{R}_{s}$ é a radiação
solar média diária e $n_{\mathrm{dias}}$ é o número de dias da janela
meteorológica.

A ação foi definida como um vetor contínuo $a \in [-1,1]^4$:
\begin{equation}
    a = [a_1, a_2, a_3, a_4].
\end{equation}
O primeiro componente define o deslocamento da data de plantio dentro da janela
de 15 de setembro a 20 de dezembro. Os demais componentes definem o manejo de
irrigação:
\begin{align}
    d_{\mathrm{plantio}} &=
    \left\lfloor
    \frac{a_1 + 1}{2} W
    \right\rceil, \\
    \theta_{\mathrm{def}} &=
    10 + \frac{a_2 + 1}{2}\cdot 80, \\
    I_{\mathrm{evento}} &=
    2 + \frac{a_3 + 1}{2}\cdot 35, \\
    I_{\mathrm{max}} &=
    20 + \frac{a_4 + 1}{2}\cdot 240,
\end{align}
em que $W$ é o tamanho da janela de plantio em dias,
$\theta_{\mathrm{def}}$ é o gatilho de déficit hídrico, $I_{\mathrm{evento}}$ é a
lâmina por evento e $I_{\mathrm{max}}$ é o limite sazonal de irrigação. A
política de irrigação foi avaliada a cada 12 dias durante 150 dias após o
plantio.

A recompensa combinou produtividade e uso de água:
\begin{equation}
    r =
    \frac{Y}{Y_{\mathrm{ref}}}
    -
    \lambda I,
\end{equation}
em que $Y$ é a produtividade corrigida da simulação DSSAT, $Y_{\mathrm{ref}} =
4200$ kg ha$^{-1}$, $I$ é a irrigação sazonal em mm e $\lambda = 0{,}0015$ é o
peso da penalização da água. Simulações com falha de execução receberam
recompensa igual a $-2$.

\subsection{Treinamento com PPO}
\label{subsec:ppo}

O agente foi treinado com Proximal Policy Optimization (PPO), usando a
implementação da biblioteca \texttt{stable-baselines3}. A política e a função de
valor foram representadas por redes neurais multicamadas com três camadas
ocultas de 128, 128 e 64 neurônios, com ativação tangente hiperbólica. A política
produz uma distribuição gaussiana sobre as quatro ações contínuas.

A função objetivo do PPO usa a razão de probabilidades entre a política nova e a
política usada na coleta das trajetórias:
\begin{equation}
    L^{\mathrm{CLIP}}(\theta) =
    \mathbb{E}_{t}
    \left[
    \min
    \left(
    \rho_t(\theta)\hat{A}_t,
    \mathrm{clip}
    \left(
    \rho_t(\theta), 1-\epsilon, 1+\epsilon
    \right)\hat{A}_t
    \right)
    \right],
\end{equation}
em que $\rho_t(\theta)$ é a razão entre probabilidades de ação,
$\hat{A}_t$ é a estimativa de vantagem e $\epsilon$ é o limite de alteração da
política em cada atualização.

Os hiperparâmetros usados no treinamento estão na
Tabela~\ref{tab:hiperparametros_ppo}.

\begin{table}[htbp]
\centering
\caption{Hiperparâmetros usados no treinamento PPO.}
\label{tab:hiperparametros_ppo}
\begin{tabular}{ll}
\hline
Parâmetro & Valor \\
\hline
Etapas totais & 500.000 \\
Ambientes paralelos & 8 \\
Passos por coleta & 256 \\
Tamanho do lote & 512 \\
Épocas por atualização & 10 \\
Taxa de aprendizado & 0,0002 \\
Fator de desconto $\gamma$ & 0,99 \\
GAE $\lambda$ & 0,95 \\
Intervalo de clip & 0,15 \\
Coeficiente de entropia & 0,01 \\
Coeficiente da função de valor & 0,50 \\
Norma máxima do gradiente & 0,50 \\
Semente aleatória & 42 \\
\hline
\end{tabular}
\end{table}

Durante o treinamento, a política foi avaliada no conjunto de validação a cada
10.000 etapas. O modelo com maior recompensa média de validação foi salvo como
melhor política. Quando a recompensa de validação não aumentava por oito
avaliações consecutivas, o treinamento era encerrado. Ao final, a política salva
foi avaliada nos conjuntos de validação e teste com ações determinísticas.

\subsection{Políticas de referência e saídas}
\label{subsec:baselines_saidas}

O desempenho do agente foi comparado com quatro políticas de referência:
\begin{enumerate}
    \item plantio em 15 de outubro sem irrigação;
    \item plantio em 15 de outubro com irrigação moderada;
    \item plantio antecipado em 25 de setembro com irrigação moderada;
    \item plantio tardio em 10 de novembro com irrigação moderada.
\end{enumerate}

As saídas computacionais foram gravadas em \texttt{outputs/}. A calibração final
foi armazenada em \texttt{outputs/calibration\_v4\_row\_bias\_corrected/},
incluindo tabelas de erro por ano, busca de $\alpha$, métricas por conjunto e
figuras em inglês. O treinamento PPO foi armazenado em
\texttt{outputs/ppo\_soja\_castro\_dssat/}, contendo modelo treinado, curva de
validação, tabelas de avaliação, comparação com políticas de referência e
figuras para discussão dos resultados.

\subsection{Reprodutibilidade}
\label{subsec:reprodutibilidade}

O fluxo final foi executado por um script único em PowerShell. A rotina primeiro
verifica se o DSSAT real executa com os arquivos de clima, solo e cultivar. Em
seguida, executa a calibração da produtividade contra o SIDRA e salva o modelo
de correção. Por fim, inicia o treinamento PPO com DSSAT real. Essa organização
mantém separados os arquivos de entrada, a calibração, o treinamento e as saídas
usadas nas tabelas e figuras do artigo.
